% SGA 3, Expose VI_B, Loop 1 English reconstruction.
% R2: independent visual review restored the intersentence
% space after the local symbolic translator's footnote.
% R3: independent PDF-link replay assigned distinct anchors
% to the eleven enumerated conclusion targets.
% Scope: Proposition 11.10bis, Remark 11.10.1,
% Proposition 11.11, and Remark 11.11.1, including
% editor's notes 119--124.
% Authority: Polo--Gille Exp6B-13oct24.pdf, component-local
% pages 87--90 (combined-reader pages 428--431; re-edition
% headers 417--420; original-margin printed pages 402--405).
% Lemma 11.12 begins on local page 90 and is excluded.
% OCR is locator-only; Reinhold e7a259f English Markdown is
% comparison-only and is credited in the component report.
% Loop 2 replaces all three source-bound commutative-diagram crops
% with faithful native LaTeX reconstructions; all are closed.

\medskip
\phantomsection\label{prop:sga3-VIB-11.10.bis}%
\noindent\textbf{命题 11.10.bis.---}%
\footnote{\label{note:sga3-VIB-119}编者注：本命题为新增，是 \textup{[Th87]}，\SGARefExternal{prop:sga3-VIB-1.4}{1.4}--\SGARefExternal{lem:sga3-VIB-1.5}{1.5} 的特例。该文作者援引了 Deligne 的一个论证\textup{（参见 \textup{[Kn71]}，III，定理 \SGARefExternal{par:sga3-VIB-1.1}{1.1}）}；也可以注意到它与 \hyperref[prop:sga3-VIB-11.8-bis]{11.8.bis} 中重述的 Serre 的论证\textup{（\textup{[Se68]}，命题 2）}的相似之处。}
\emph{设 \(S\) 为 Noether 概形，\(G\) 为在 \(S\) 上平坦、拟紧且拟分离的 \(S\)-群，\(\mathcal E\) 为拟凝聚 \(G\)-\(\mathcal O_S\)-模，\(\mathcal M\) 为 \(\mathcal E\) 的凝聚 \(\mathcal O_S\)-子模。则 \(\mathcal M\) 包含于一个在 \(G\) 下稳定的凝聚 \(\mathcal O_S\)-子模中。}

\emph{证明。}
用 \(f\) 表示态射 \(G\to S\)，用 \(\tau\) 表示伴随态射
\[
  \mathcal E\longrightarrow f_*f^*(\mathcal E).
\]
由 \hyperref[par:sga3-VIB-11.6]{11.6}，\(\mathcal E\) 上的 \(G\)-\(\mathcal O_S\)-模结构由 \(\mathcal O_G\)-模 \(f^*(\mathcal E)\) 的一个自同构 \(\theta\) 给出，使得态射
\[
  \delta=\theta\circ\tau
\]
满足 \hyperref[par:sga3-VIB-11.6]{11.6} 的条件\hyperref[diag:sga3-VIB-11.6-1]{\textup{(1)}}和\hyperref[diag:sga3-VIB-11.6-2]{\textup{(2)}}。（\textup{[Th87]} 考虑的情形更一般：该文作者考虑一个 \(G\)-概形 \(X\) 和一个 \(G\)-等变 \(\mathcal O_X\)-模 \(\mathcal E\)\textup{（参见\SGARefLink{sga3:i:section:6}{第 I 篇，第 6 节}）}；这里 \(X=S\)，带有 \(G\) 的平凡作用。）

由于 \(S\) 是 Noether 概形，\(\mathcal E\) 是其凝聚子模 \(\mathcal F_\alpha\) 的滤过归纳极限\textup{（参见 EGA~I，9.4.9）}。于是 \(f^*(\mathcal E)\) 是诸 \(f^*(\mathcal F_\alpha)\) 的滤过归纳极限；由于 \(f\) 平坦，后者都是 \(f^*(\mathcal E)\) 的子模。此外，由于 \(f\) 拟紧且拟分离，由 \hyperref[par:sga3-VIB-11.0]{11.0}，\(\mathcal O_S\)-模 \(f_*f^*(\mathcal E)\) 拟凝聚，并且是拟凝聚子模 \(f_*f^*(\mathcal F_\alpha)\) 的滤过归纳极限。因此，\(\mathcal E\) 是拟凝聚 \(\mathcal O_S\)-子模 \(\mathcal E_\alpha\) 的滤过归纳极限，其中 \(\mathcal E_\alpha\) 表示 \(f_*f^*(\mathcal F_\alpha)\) 在 \(\delta\) 下的逆像；等价地，它是下列复合的核：
\phantomsection\label{form:sga3-VIB-11.10bis-delta-alpha}%
\[
  \overline{\delta}_\alpha:\quad
  \mathcal E\xrightarrow{\ \delta\ }f_*f^*(\mathcal E)
    \longrightarrow f_*f^*(\mathcal E/\mathcal F_\alpha).
\]
由于 \(\mathcal M\) 凝聚且 \(S\) 为 Noether 概形，\(\mathcal M\) 的每个递增子模列都趋于稳定；所以 \(\mathcal M\) 包含于某个 \(\mathcal E_\alpha\) 中。我们将证明每个 \(\mathcal E_\alpha\) 都凝聚且在 \(G\) 下稳定。

设
\[
  u:f^*(\mathcal E)\longrightarrow\varepsilon_*(\mathcal E)
\]
为经伴随对应于从 \(\mathcal E\) 到 \(f_*\varepsilon_*(\mathcal E)=\mathcal E\) 的恒等态射的态射。则复合
\phantomsection\label{form:sga3-VIB-11.10bis-adjunction-retract}%
\[
  \mathcal E\xrightarrow{\ \delta\ }f_*f^*(\mathcal E)
    \xrightarrow{\ f_*(u)\ }f_*\varepsilon_*(\mathcal E)
    =\mathcal E
\]
为恒等态射%
\footnote{\label{source-note:sga3-VIB-11.10bis-delta-arrows}源 PDF 将此复合中的第一个箭头及下图中的下方水平箭头标为 \(\tau\)。这两处必须为 \(\delta=\theta\circ\tau\)：条件 11.6\textup{(2)} 给出 \(f_*(u)\circ\delta=\operatorname{id}_{\mathcal E}\)，而 \(\mathcal E_\alpha\) 的定义保证方块交换。}
\textup{（参见 \hyperref[diag:sga3-VIB-11.6-2]{11.6\textup{(2)}}）}。有如下交换图
\begin{center}
  \phantomsection\label{diag:sga3-VIB-11.10bis-coherence}%
  \begin{tikzcd}[row sep=large,column sep=large]
    \mathcal E_\alpha
      \arrow[r,"\delta"]
      \arrow[d,"i_\alpha"'] &
    f_*f^*(\mathcal F_\alpha)
      \arrow[r,"f_*(u)"]
      \arrow[d] &
    \mathcal F_\alpha
      \arrow[d,"j_\alpha"] \\
    \mathcal E
      \arrow[r,"\delta"'] &
    f_*f^*(\mathcal E)
      \arrow[r,"f_*(u)"'] &
    \mathcal E\mathpunct{,}
  \end{tikzcd}
\end{center}
其中 \(i_\alpha\)\textup{（相应地，\(j_\alpha\)）}表示 \(\mathcal E_\alpha\)\textup{（相应地，\(\mathcal F_\alpha\)）}到 \(\mathcal E\) 中的包含态射。由此可知 \(i_\alpha\) 通过 \(j_\alpha\) 分解；因此 \(\mathcal E_\alpha\) 是 \(\mathcal F_\alpha\) 的子模，从而凝聚。

最后证明 \(\mathcal E_\alpha\) 在 \(G\) 下稳定\textup{（参见 \hyperref[rem:sga3-VIB-11.7.bis]{11.7.bis}）}。用 \(h:G\to S\) 表示 \(f:G\to S\) 的第二个副本，并考虑下列交换图：
\begin{center}
  \phantomsection\label{diag:sga3-VIB-11.10bis-projections}%
  \begin{tikzcd}[row sep=large,column sep=large]
    G\times_S G
      \arrow[r,"q"]
      \arrow[d,"p"']
      \arrow[dr,"\varphi" description] &
    G \arrow[d,"h"] \\
    G \arrow[r,"f"'] &
    S\mathpunct{.}
  \end{tikzcd}
\end{center}
正合列
\phantomsection\label{seq:sga3-VIB-11.10bis-first-exact}%
\[
  0\longrightarrow\mathcal E_\alpha
    \longrightarrow\mathcal E
    \xrightarrow{\ \overline{\delta}_\alpha\ }
    h_*h^*(\mathcal E/\mathcal F_\alpha)
\]
由于 \(f\) 平坦，给出正合列
\phantomsection\label{seq:sga3-VIB-11.10bis-dagger}%
\[
  \tag{\(\dagger\)}
  0\longrightarrow f_*f^*(\mathcal E_\alpha)
    \longrightarrow f_*f^*(\mathcal E)
    \xrightarrow{\ f_*f^*(\overline{\delta}_\alpha)\ }
    f_*f^*h_*h^*(\mathcal E/\mathcal F_\alpha).
\]
此外，由于 \(h\) 拟紧且拟分离，而 \(f\) 平坦，典范态射
\phantomsection\label{form:sga3-VIB-11.10bis-flat-base-change}%
\[
  f^*h_*h^*(\mathcal E/\mathcal F_\alpha)
    \longrightarrow
  p_*q^*h^*(\mathcal E/\mathcal F_\alpha)
    =p_*\varphi^*(\mathcal E/\mathcal F_\alpha)
\]
是同构。因此，\hyperref[seq:sga3-VIB-11.10bis-dagger]{\textup{(\(\dagger\))}}的最右项是 \(\varphi_*\varphi^*(\mathcal E/\mathcal F_\alpha)\)。采用 \hyperref[par:sga3-VIB-11.6]{11.6} 的记号，并用 \(\pi_\alpha\) 表示投影 \(\mathcal E\to\mathcal E/\mathcal F_\alpha\)，可得如下交换图，其下行正合：
\begin{center}
  \phantomsection\label{diag:sga3-VIB-11.10bis-stability}%
  {\small
  \begin{tikzcd}[row sep=large,column sep=large]
    {} &
    \mathcal E_\alpha
      \arrow[r,"\delta"]
      \arrow[d,"\delta"'] &
    f_*f^*(\mathcal E)
      \arrow[d,"{f_*(f^*(\pi_\alpha)\boxtimes\delta_G)}"] \\
    f_*f^*(\mathcal E_\alpha)
      \arrow[r] &
    f_*f^*(\mathcal E)
      \arrow[r,"{f_*f^*(\overline{\delta}_\alpha)}"'] &
    \varphi_*\varphi^*(\mathcal E/\mathcal F_\alpha)
  \end{tikzcd}}
\end{center}
由于
\[
  f_*\bigl(f^*(\pi_\alpha)\boxtimes\delta_G\bigr)\circ\delta
\]
在 \(\mathcal E_\alpha\) 上为零，可知 \(\delta\) 将 \(\mathcal E_\alpha\) 映入 \(f_*f^*(\mathcal E_\alpha)\)；换言之，\(\mathcal E_\alpha\) 在 \(G\) 下稳定。命题得证。

\medskip
\phantomsection\label{rem:sga3-VIB-11.10.1}%
\noindent\textbf{注 11.10.1.---}%
\footnote{\label{note:sga3-VIB-120}编者注：原\hyperref[rem:sga3-VIB-11.10.1]{注 11.10.1} 的第一部分已纳入 \hyperref[lem:sga3-VIB-11.8]{11.8} 和 \hyperref[prop:sga3-VIB-11.9]{11.9}\textup{（将``自由''改为``投射''）}；下面的反例勘正了其第二部分。}
在 \hyperref[lem:sga3-VIB-11.8]{11.8} 中，仅假设 \(C\) 为平坦 \(A\)-模并不够，即使 \(A\) 是主理想整环也如此。事实上，O.\ Gabber 和 J.-P.\ Serre 分别向我们指出了下列反例。

\smallskip
\phantomsection\label{rem:sga3-VIB-11.10.1-a}%
\noindent\textup{(a)}
设 \((A,\mathfrak m)\) 为离散赋值环，\(K\) 为其分式域，\(G\) 为由 \(K\)-群 \((\mathbb Z/2\mathbb Z)_K\) 经``补零延拓''得到的 \(A\)-群。常值群 \((\mathbb Z/2\mathbb Z)_A\)，因而其子群 \(G\)，通过交换 \(v_1\) 与 \(v_2\) 作用于以 \(v_1,v_2\) 为基的自由 \(A\)-模 \(V\)。于是 \(Av_1\) 不是 \(V\) 的 \(G\)-子模，但它是下列 \(G\)-子模的交：
\[
  Av_1+\mathfrak m^n v_2,\qquad n\in\mathbb N^*.
\]
因此，不存在 \(V\) 中包含 \(v_1\) 的最小 \(G\)-子模。

\smallskip
\phantomsection\label{rem:sga3-VIB-11.10.1-b}%
\noindent\textup{(b)}
设 \(A\) 为异于其分式域 \(K\) 的整环，\(G\) 为对应于下列 Hopf 代数的平坦仿射 \(A\)-群：
\phantomsection\label{form:sga3-VIB-11.10.1-Hopf-example}%
\[
  \mathcal A(G)=\{P\in K[T]\mid P(0)\in A\}.
\]
其余乘法、余单位和对极分别定义为
\[
  \Delta(T)=T\otimes1+1\otimes T,\qquad
  \varepsilon(T)=0,\qquad
  \tau(T)=-T.
\]
\textup{（因此 \(G\otimes_AK=\mathbb G_{a,K}\)。）}

设 \(V=Av_1\oplus Av_2\) 在 \(A\) 上自由，\(u\) 为 \(V\) 的自同态，定义为
\[
  u(v_1)=v_2,\qquad u(v_2)=0,
\]
于是 \(u^2=0\)。此时 \(V\) 具有下列 \(\mathcal A(G)\)-余模结构：
\phantomsection\label{form:sga3-VIB-11.10.1-coaction}%
\[
  \mu(m)=m\otimes1+u(m)\otimes T.
\]
\(V\) 中包含 \(v_1\) 的 \(G\)-子模恰好是形如 \(Av_1\oplus Iv_2\) 的 \(A\)-子模，其中 \(I\) 为 \(A\) 的非零理想。它们的交为 \(Av_1\)，这不是 \(G\)-子模。因此，不存在 \(V\) 中包含 \(v_1\) 的最小 \(G\)-子模。此外，
\[
  C=A\oplus K\cdot T
\]
是 \(\mathcal A(G)\) 的子余代数，在 \(A\) 上平坦，且余作用
\[
  \mu:V\longrightarrow V\otimes\mathcal A(G)
\]
通过 \(V\otimes C\) 分解。因此，这个``极其简单''的余代数 \(C\) 也给出一个反例。

最后，以上两个例子都是下列构造的特例。设 \(A\) 为异于其分式域 \(K\) 的整环，\(B\) 为在 \(A\) 上自由的 Hopf \(A\)-代数。记
\[
  \varepsilon:B\longrightarrow A,\qquad
  I=\ker(\varepsilon)
\]
为增广及增广理想。由于 \(B=A\cdot1\oplus I\)，容易看出
\phantomsection\label{form:sga3-VIB-11.10.1-B-prime}%
\[
  B'=\{b\in B\otimes_AK\mid \varepsilon(b)\in A\}
\]
是 \(B_K\) 的 Hopf 子代数。假设 \(V\) 为一个 \(B\)-余模，作为 \(A\)-模是以 \((v_1,\ldots,v_n)\) 为基的自由模，且 \(\mu_V(v_1)\neq v_1\otimes1\)。则 \(Av_1\) 不是 \(V\) 的子余模，但它是下列子余模的交：
\[
  Av_1+J\cdot V
\]
其中 \(J\) 遍历 \(A\) 的非零理想。%
\begingroup
\renewcommand{\thefootnote}{\fnsymbol{footnote}}%
\footnote[1]{\label{note:sga3-VIB-11.10.1-ideal-letter}
译者注：法文重版用 \(I\) 同时表示上面固定的增广理想和最后这个交中所遍历的理想。我们将后者记作 \(J\)，以区分这两种用途。}%
\endgroup\ 
因此，不存在 \(V\) 中包含 \(v_1\) 的最小子余模。

\medskip
\phantomsection\label{prop:sga3-VIB-11.11}%
\noindent\textbf{命题 11.11.---}
\emph{设 \(G\) 为域 \(k\) 上的代数群。下列条件等价：}
\begin{enumerate}
\item[\textup{(i)}]\phantomsection\label{concl:sga3-VIB-11.11-i}
  \emph{\(G\) 仿射。}
\item[\textup{(ii)}]\phantomsection\label{concl:sga3-VIB-11.11-ii}
  \emph{\(G\) 拟仿射。}
\item[\textup{(iii)}]\phantomsection\label{concl:sga3-VIB-11.11-iii}
  \emph{\(G\) 忠实地作用于一个拟仿射 \(k\)-概形 \(X\)。}
\item[\textup{(iv)}]\phantomsection\label{concl:sga3-VIB-11.11-iv}
  \emph{\(G\) 线性且忠实地作用于一个 \(k\)-向量空间\textup{（不必有限维）}。}
\item[\textup{(v)}]\phantomsection\label{concl:sga3-VIB-11.11-v}
  \emph{\(G\) 同构于 \(\operatorname{GL}(n)_k\) 的一个闭子群。}
\end{enumerate}

\emph{证明。}
蕴含关系\hyperref[concl:sga3-VIB-11.11-i]{\textup{(i)}}\(\Rightarrow\)\hyperref[concl:sga3-VIB-11.11-ii]{\textup{(ii)}}显然，而\hyperref[concl:sga3-VIB-11.11-ii]{\textup{(ii)}}\(\Rightarrow\)\hyperref[concl:sga3-VIB-11.11-iii]{\textup{(iii)}}是因为 \(G\) 通过平移忠实地作用于自身。

假设 \(G\) 忠实地右作用于一个拟仿射 \(k\)-概形 \(X\)。%
\footnote{\label{note:sga3-VIB-121}编者注：以下对原文作了扩充。我们还选择了 \(G\) 在 \(X\) 上的右作用，以便得到 \(G\) 在 \(\mathcal O(X)\) 上的左线性作用。}
由于 \(X\) 拟仿射，它分离且拟紧。%
\footnote{\label{note:sga3-VIB-122}编者注：回忆，一个 \(k\)-概形 \(X\) 称为拟仿射的，是指它同构于一个仿射 \(k\)-概形的拟紧开子概形\textup{（EGA~II，5.1.1）}。}
同样，\(G\) 分离\textup{（\hyperref[sga3:VIA:0.3]{VI\(_A\)，0.3}）}且拟紧\textup{（因为它在 \(k\) 上有限型）}。所以由\hyperref[concl:sga3-VIB-11.1-c]{11.1\textup{(c)}}，有典范同构
\phantomsection\label{form:sga3-VIB-11.11-functions-products}%
\[
  \mathcal O(X\times G)
    =\mathcal O(X)\otimes\mathcal O(G),
  \qquad
  \mathcal O(X\times G\times G)
    =\mathcal O(X)\otimes\mathcal O(G)\otimes\mathcal O(G).
\]
由此可知，态射
\[
  \mu:\mathcal O(X)\longrightarrow
    \mathcal O(X)\otimes\mathcal O(G)
\]
（由 \(X\times G\to X\) 诱导）赋予 \(\mathcal O(X)\) 一个右 \(\mathcal O(G)\)-余模结构；即 \(G\) 左线性作用于 \(k\)-代数 \(\mathcal O(X)\)。因此，\(G\) 也右作用于仿射包络
\[
  X_{\mathrm{af}}=\operatorname{Spec}\mathcal O(X)
\]
（这是 \(X\) 的仿射包络），且典范态射 \(X\to X_{\mathrm{af}}\) 是 \(G\)-等变的。

由于 \(X\) 拟仿射，态射 \(X\to X_{\mathrm{af}}\) 是开浸入\textup{（EGA~II，5.1.2）}。所以 \(G\) 在 \(X_{\mathrm{af}}\) 上的作用当然也是忠实的，从而 \(G\) 在 \(k\)-代数 \(\mathcal O(X)=\mathcal O(X_{\mathrm{af}})\) 上的左线性作用忠实。这就证明了
\hyperref[concl:sga3-VIB-11.11-iii]{\textup{(iii)}}
\(\Rightarrow\)
\hyperref[concl:sga3-VIB-11.11-iv]{\textup{(iv)}}.

现在假设 \(G\) 忠实地作用于一个 \(k\)-向量空间 \(V\)。由 \hyperref[cor:sga3-VIB-11.10]{11.10}，\(V\) 是在 \(G\) 下稳定的有限维向量子空间 \(V_i\) 的归纳极限。设 \(K_i\) 为 \(G\) 在 \(V_i\) 上诱导的作用的核；即 \(G\to\operatorname{Aut}(V_i)\) 的核。则 \(K_i\) 是 \(G\) 的闭子概形，而 \(G\) 的作用忠实意味着诸 \(K_i\) 的交是 \(G\) 的单位子群。由于 \(G\) 为 Noether 概形，已有某个 \(K_i\) 是单位子群。因此 \(G\to\operatorname{Aut}(V_i)\) 是单态射，从而由 \hyperref[cor:sga3-VIB-1.4.2]{1.4.2}，它是闭浸入。这就证明了
\hyperref[concl:sga3-VIB-11.11-iv]{\textup{(iv)}}
\(\Rightarrow\)
\hyperref[concl:sga3-VIB-11.11-v]{\textup{(v)}}。最后，
\hyperref[concl:sga3-VIB-11.11-v]{\textup{(v)}}
\(\Rightarrow\)
\hyperref[concl:sga3-VIB-11.11-i]{\textup{(i)}} 显然成立，命题得证。

\medskip
\phantomsection\label{rem:sga3-VIB-11.11.1}%
\noindent\textbf{注 11.11.1.---}
命题 \hyperref[prop:sga3-VIB-11.11]{11.11} 有如下推广。设 \(S\) 为维数至多为 \(1\) 的正则且局部 Noether 的概形，\(G\) 为在 \(S\) 上平坦、拟紧且拟分离的群概形。\textup{（此时，\(\mathcal A(G)\) 是无挠 \(\mathcal O_S\)-模，因而平坦。）}

\begin{enumerate}
\item[\textup{(a)}]\phantomsection\label{concl:sga3-VIB-11.11.1-a}
  下列条件等价：%
  \footnote{\label{note:sga3-VIB-123}编者注：这些条件的等价性在新增的\hyperref[sec:sga3-VIB-12]{第 12 节}中证明。}
  \begin{enumerate}
  \item[\textup{(i)}]\phantomsection\label{concl:sga3-VIB-11.11.1-a-i}
    \(G\) 在 \(S\) 上仿射。
  \item[\textup{(ii)}]\phantomsection\label{concl:sga3-VIB-11.11.1-a-ii}
    \(G\) 在 \(S\) 上拟仿射。
  \item[\textup{(iii)}]\phantomsection\label{concl:sga3-VIB-11.11.1-a-iii}
    \(G\) 忠实地作用于一个拟仿射平坦 \(S\)-概形。
  \item[\textup{(iv)}]\phantomsection\label{concl:sga3-VIB-11.11.1-a-iv}
    \(G\) 线性且忠实地作用于一个平坦拟凝聚 \(\mathcal O_S\)-模。
  \end{enumerate}

\item[\textup{(b)}]\phantomsection\label{concl:sga3-VIB-11.11.1-b}
  如果 \(G\) 还在 \(S\) 上有限型，且 \(S\) 为 Noether 概形，则这些条件蕴含 \(G\) 同构于 \(\operatorname{Aut}(\mathcal E)\) 的一个闭子群，其中 \(\mathcal E\) 为有限秩局部自由 \(\mathcal O_S\)-模。%
  \footnote{\label{note:sga3-VIB-124}编者注：这一结果及其若干推广在新增的\hyperref[sec:sga3-VIB-13]{第 13 节}中证明。}
\end{enumerate}
